%% 第六章

\chapter{总结与展望}
\chaptermark{总结与展望}

\section{研究工作总结}

本文围绕深空场景的三维重建与新视角合成任务展开研究，重点关注3D Gaussian Splatting在火星等深空场景中的适用性问题。深空图像通常具有远近尺度跨度大、纹理重复且局部差异弱、跨视角光照条件不稳定等特点，这些因素会改变3DGS优化过程中的监督分布，使远区域细节、显式高斯形态和颜色一致性更容易出现退化。基于此，本文从数据特征和3DGS优化机制出发，将深空场景中的主要问题归纳为远区域监督不足、高斯基元形态不稳定和跨视角外观差异三个方面。

针对上述问题，本文提出DSS-GS深空场景重建与新视角合成框架。该框架以3D Gaussian Splatting为主干，保留其显式三维高斯表示和快速可微渲染优势，同时在训练阶段引入面向深空场景数据的约束。其中，FAME远区域增强模块基于辅助深度图的相对排序构造远区域掩膜$M_{\mathrm{far}}$，并在基础重建损失之外单独计算远区域颜色误差，使远处地表、弱纹理结构和背景层次获得更明确的监督信号；RVS高斯形态与尺度稳定模块由有效秩正则和体积正则组成，分别约束高斯尺度方向分布和整体空间覆盖范围，用于减少异常拉伸、过度膨胀和局部尺度突变。对于存在明显光照扰动的情况，本文还设计了跨视角外观一致性建模流程，通过参考图像特征和动态外观表示缓解颜色偏移与亮度变化带来的影响。

在实验验证方面，本文基于M3arsSynth数据集划分构建实验流程，共使用593个视频场景，其中416个场景用于新视角合成质量对比和模块消融实验，177个场景用于构造光照扰动数据并评估外观一致性分支。实验过程中，各方法使用相同的多视角图像、COLMAP相机参数、初始稀疏点云、训练测试划分和评价指标，以保证结果具有可比性。评价内容覆盖三维场景表示质量、新视角合成质量、光照扰动鲁棒性和模块消融等方面。

实验结果表明，本文方法在深空场景新视角合成任务上取得了较好的综合表现。在整体质量对比中，DSS-GS在$\mathrm{SSIM}$、$\mathrm{PSNR}$和$\mathrm{LPIPS}$三项指标上分别达到0.961、38.49dB和0.078，优于3DGS及其他对比方法；在三维表示分析中，重投影深度误差由3DGS的0.287降低至0.242，说明空间层次和深度投影一致性有所改善。难度分组实验表明，本文方法在困难组、常规组和简单组中均取得正向收益，其中困难组提升最明显。光照扰动实验和消融实验进一步说明，外观一致性建模能够提升光照扰动条件下的渲染稳定性，FAME与RVS则分别从远区域监督增强和高斯形态稳定两个角度发挥作用，二者联合使用时取得最优综合性能。

综上，本文将深空场景的数据特征与3DGS优化机制联系起来，围绕远区域监督、三维表示稳定和外观一致性三个方向设计DSS-GS方法，并通过定量指标、可视化结果、重投影一致性和消融实验验证了所设计模块的有效性。本文工作为深空场景的三维重建和自由视角观察提供了一种基于显式高斯表示的改进方案。

\section{研究前景展望}

本文方法仍存在进一步完善空间。首先，FAME远区域增强模块依赖辅助深度图构造远区域掩膜$M_{\mathrm{far}}$。虽然基于相对深度排序的方式能够避免固定绝对深度阈值带来的跨场景适应性问题，但当单目深度估计在低纹理、遮挡或阴影区域出现偏差时，远区域划分的可靠性仍可能下降。后续可结合图像纹理强度、稀疏点云密度、相机视角分布和多视图重投影关系共同判断远区域范围，并引入置信度加权或动态比例策略，使远区域监督更加稳定。

其次，RVS主要从高斯尺度参数出发约束显式表示的形态和空间覆盖范围。有效秩正则和体积正则能够减少异常拉伸和局部尺度突变，但仍不能完全替代真实几何监督或显式多视图一致性约束。未来可以进一步引入跨视角深度一致性、局部表面平滑约束、法向一致性或稀疏点云置信度约束，使高斯形态稳定受到更直接的几何关系引导，从而进一步减少局部空洞、边界断裂和深度层次波动。

再次，跨视角外观一致性建模仍有继续改进空间。深空场景中的外观变化可能同时包含曝光变化、阴影分布、白平衡偏移、材质反射差异和成像设备响应差异，当前流程仍难以完全区分场景固有颜色与外部成像条件。后续可利用连续视角或视频序列中的外观变化规律，引入更细致的亮度、阴影和材质建模，使模型在复杂光照条件下保持更稳定的颜色表达和局部细节。

此外，本文实验主要围绕M3arsSynth火星场景多视角序列展开，后续还可以在更多类型的深空场景数据上进行验证，包括不同地貌尺度、不同相机轨迹、不同光照条件和更复杂遮挡关系的数据。通过扩展数据来源和评价维度，可以更全面地分析DSS-GS在行星、月面地形或小天体等场景中的适用范围，并进一步检验远区域增强、形态稳定和外观建模之间的互补关系。
